\section{Are the Features Useful for Planning?}
\label{sec:planning}

The energy score measures how well the probe predicts the logged future. A
planner is judged differently: a trajectory that differs from the driver's but
stays safe and makes progress is fine. We therefore scored the probes with EPDMS.

\subsection{The diffusion probe on navtest}
\label{sec:epdms-probe}

\begin{table}[htbp]
\centering
\caption{EPDMS of the diffusion probes on navtest (12,146 scenes, official
devkit). ``Sample'' scores the first of 32 samples, ``mean'' their average,
``best of 4'' the best of the first four per scene (an oracle). Sub-scores are
for the mean.}
\label{tab:epdms-probe}
\begin{tblr}{colspec={lrrrrrrr}, row{1}={font=\bfseries}, cells={font=\small}}
\toprule
Input & Sample & Mean & Best of 4 & NC & DAC & EP & EC \\
\midrule
constant velocity           & 0.259 &       &       & 0.681 & 0.578 & 0.777 & 0.553 \\
none                        & 0.316 & 0.367 &       & 0.742 & 0.693 & 0.787 & 0.286 \\
ego status + command        & 0.578 & 0.645 & 0.800 & 0.923 & 0.765 & 0.861 & 0.848 \\
V-JEPA 2                    & 0.618 & 0.694 &       & 0.921 & 0.848 & 0.867 & 0.607 \\
$q_{\text{sem}}$            & 0.620 & 0.709 &       & 0.933 & 0.871 & 0.862 & 0.515 \\
DINOv2                      & 0.620 & 0.710 &       & 0.928 & 0.874 & 0.871 & 0.517 \\
WM23                        & 0.662 & 0.740 &       & 0.941 & 0.870 & 0.877 & 0.689 \\
WM11-ctx                    & 0.722 & 0.788 & 0.882 & 0.963 & 0.890 & 0.875 & 0.788 \\
\midrule
$q_{\text{sem}}$ + status + command & 0.711 & 0.775 & 0.877 & 0.952 & 0.879 & 0.874 & 0.877 \\
WM11-ctx + status + command  & \textbf{0.756} & \textbf{0.810} & \textbf{0.900} & 0.969 & 0.898 & 0.876 & 0.877 \\
\bottomrule
\end{tblr}
\end{table}

The image-only probes in Table~\ref{tab:epdms-probe} rank the same way on EPDMS as on the energy score, with the
world model clearly first (0.788 against 0.710 for DINOv2). The offline ranking
of Section~\ref{sec:features} is not an artefact of the metric.

The ego-status probe, which had the best energy score, scores only 0.645. It
leaves the drivable area much more often (DAC 0.765 against 0.890) because it
cannot see the road. So its advantage on the energy score comes from predicting
speed well, and speed is not what makes a plan safe. This is the clearest
evidence in this report that displacement-type metrics can rank planners wrongly.

The mean of the samples beats a single sample by 0.04 to 0.09. Most of this is
extended comfort: samples are drawn independently in each frame, so consecutive
plans disagree. A sampler that is right about the uncertainty is penalized for
showing it.

The best of four samples reaches 0.882 for WM11-ctx and 0.900 with the ego
status. In most scenes one of the samples is a good plan. The weak step is
choosing it.

Adding the ego status and command to the world-model probe lowers the energy
score by 41\%, but raises EPDMS only from 0.788 to 0.810, and mostly through
comfort (EC 0.788 to 0.877). The speed input changes smoothly from frame to frame
and ties consecutive plans to the same speed. DAC barely moves, because the probe
already knew the speed from the image, and its remaining failures are about the
road. With the ego status, the world-model features are still worth 0.035 over
$q_{\text{sem}}$.

\paragraph{navhard.}
On the two-stage navhard split all probes score between 0.33 and 0.41 with the
sample mean, and the image-only world-model probe (0.367) is only a little above
constant velocity under the same devkit (0.326). With the ego status both encoders
reach about 0.41. The loss happens in stage two, where lane keeping falls to
about 0.45 for every planner, constant velocity included. Stage two starts from a
displaced ego pose, and a probe that imitates the logged driver has never seen
how to return to the lane.

\subsection{A fast planner on the frozen features}
\label{sec:fast-planner}

The diffusion probe was designed to study distributions, not to drive. Its
weaknesses as a planner are the ones above: independent samples per frame and no
way to pick the best sample. To measure how much driving quality the frozen
features support, we built a planner for that purpose. We also wanted it to train
in minutes, so that many variants can be compared.

\paragraph{Fast EPDMS.}
A full navtest evaluation takes about 20 minutes, too long to use during
training. We wrote a scorer that loads the metric cache of a fixed set of scenes
once and scores all candidate trajectories of a scene in one call. It reproduces
the official per-scene scores to four decimals. To choose the subset size, we
compared the subset mean with the full navtest score for about 20 planners we had
already scored, over 50 random draws per size. With 500 scenes the mean absolute
error is 0.017 and the Spearman rank correlation with full navtest is 0.98, while
with 25 scenes the error is 0.06, larger than most differences between our runs.
Scoring 500 scenes takes about one minute. Final numbers are always on full
navtest.

\paragraph{Planner.}
The ego velocity, acceleration, past poses and driving command go through an MLP
into one ego token. Each cell of the cached feature grid becomes one image token.
Two transformer layers mix them. A regression head reads one query and outputs
the eight poses, trained with an $L_1$ loss. A vocabulary head, in the style of
VADv2 and Hydra-MDP~\cite{vadv2,hydra_mdp}, scores 256 anchor trajectories from
k-means on the training futures and refines the best one. All 92k training scenes
fit on the GPU, and one run of 6000 steps takes one to four minutes.

\begin{table}[htbp]
\centering
\caption{Fast planner, EPDMS on the 500-scene navtest subset. ``Best of 8'' is the
best of the top 8 anchors of the vocabulary head (an oracle). The first block
gives the reference points scored on the same scenes.}
\label{tab:fast-planner}
\begin{tblr}{colspec={lrr}, row{1}={font=\bfseries}}
\toprule
Planner & EPDMS & Best of 8 \\
\midrule
human (logged future)                         & 0.949 & \\
PDM-Closed (privileged map and agents)        & 0.871 & \\
diffusion probe WM11-ctx, sample mean         & 0.782 & \\
constant velocity                             & 0.209 & \\
\midrule
regression, WM11-ctx, width 256               & \textbf{0.887} & \\
vocabulary, WM11 (3 seeds)                    & 0.863--0.885 & 0.956 \\
vocabulary, WM11-ctx                          & 0.853 & 0.949 \\
\quad without velocity and acceleration       & 0.831 & 0.946 \\
\quad with WM23-ctx                           & 0.836 & 0.944 \\
\quad with DINOv2                             & 0.828 & 0.944 \\
\quad with $q_{\text{sem}}$                   & 0.824 & 0.955 \\
\quad with V-JEPA 2                           & 0.821 & 0.942 \\
\quad without image                           & 0.720 & 0.909 \\
regression, WM11-ctx, width 128              & 0.869 & \\
\quad with 10k instead of 92k training scenes & 0.824 & \\
\bottomrule
\end{tblr}
\end{table}

\paragraph{Results.}
Table~\ref{tab:fast-planner} gives the main runs. On full navtest the
regression planner reaches 0.881 after 1.8 minutes of training, and the
vocabulary planner 0.857. Both are well above the diffusion probe (0.788 on full
navtest), and on the subset above the privileged PDM-Closed (0.871).

The ablations repeat the earlier findings in the planning metric. Without image
features EPDMS falls to about 0.72, so the image is needed for driving even with
the full ego status. Without velocity and acceleration it falls by 0.02. The world
model is again the best encoder, ahead of DINOv2, $q_{\text{sem}}$ and V-JEPA~2.
Training data matters more than model details: 10k instead of 92k scenes costs
0.045, while 20k instead of 6k steps gains 0.01, the vocabulary size between 64 and
1024 changes little, and the regression and vocabulary heads are equal within
seed noise (three seeds span 0.863 to 0.885).

The best of the top 8 anchors scores 0.95 to 0.96, the level of the human driver.
As with the diffusion probe, the proposals contain a good trajectory and imitation
is not enough to rank them. Hydra-MDP addresses this by learning the PDM
sub-scores of every anchor. We tried this with sub-score targets for 5000 training
scenes (5\% of the data, which took about 50 minutes to compute). On these scenes
the best anchor per scene scores 0.991 and the best single anchor for all scenes
only 0.661, so the right choice depends on the scene. But distillation gave
0.872 to 0.877, within the seed noise of the plain vocabulary planner. With 5\% of
the targets the sub-score head does not learn enough to rank better than
imitation. Targets for all of navtrain would take about 15 hours to compute, and
we stopped there.

\subsection{Comparison with published planners}
\label{sec:published}

\begin{table}[htbp]
\centering
\caption{navtest EPDMS. Published numbers from the respective papers. Our devkit
scores the human and PDM-Closed references a few points higher than published
(Section~\ref{sec:epdms}), so the comparison is approximate.}
\label{tab:published}
\begin{tblr}{colspec={lllr}, row{1}={font=\bfseries}, cells={font=\small}}
\toprule
Planner & Sensors & Trained & EPDMS \\
\midrule
ego-status MLP                         & none           & end to end & 0.640 \\
TransFuser~\cite{transfuser}           & 3 cameras + LiDAR & end to end & 0.767 \\
Hydra-MDP++~\cite{hydra}               & cameras        & end to end & 0.814 \\
PDM-Closed~\cite{pdm}                  & privileged     & rule-based & 0.828 \\
DiffusionDrive~\cite{diffusiondrive}   & cameras + LiDAR & end to end & 0.843 \\
human                                  &                &            & 0.903 \\
\midrule
ours, diffusion probe, WM11-ctx        & front camera   & frozen + 19k-param probe & 0.788 \\
ours, diffusion probe, WM11-ctx + status & front camera & frozen + probe & 0.810 \\
ours, fast regression planner          & front camera   & frozen + 2 min training & 0.881 \\
\bottomrule
\end{tblr}
\end{table}

Table~\ref{tab:published} places our planners next to published ones. The
published planners train their own backbones end to end, most with more cameras
or LiDAR. Our planners see only the front camera through a frozen model that was
never trained on trajectories. Even with a few points of devkit difference, the
fast planner on Orbis features is in the range of strong end-to-end planners.
This is the practical answer to our question: the world model has learned enough
about driving from video alone that a two-minute read-out drives as well as
planners trained for the task.
